\documentclass{article}
\usepackage{pseudocode}
\begin{document}

\section{TAD Code Binaire}
\begin{tad}
     \tadNom{Code Binaire}
        
     \tadDependances{TableDeCodage, Caractere }
     \begin{tadOperations}{obtenirCodeBinaire}
        \tadOperation{obtenirCodeBinaire}{\tadDeuxParams{TableDeCodage}{Caractere}}{\tadParam{CodeBinaire}}
        \tadOperationAvecPreconditions{estBinaire}{\tadParam{CodeBinaire}}{\tadParam{Booléen}}
     \end{tadOperations}
      \begin{tadSemantiques}{}
            obtenirCodeBinaire : Obtient grâce à la table de codage le code binaire d'un caractère 
    estBinaire : Permet de vérifier que le code est bien binaire
         
      \end{tadSemantiques}
    \begin{tadPreconditions}{obtenirFilsGauche(A)}
            \tadPrecondition{estBinaire(code)}{non estVide(code)}
        \end{tadPreconditions}

\end{tad}   


\section{TAD Octet}
\begin{tad}
\tadNom{Octet}
\tadDependances{Chaine de caracteres, Code Binaire, Booleen}
\begin{tadOperations}{estOctet}
    \tadOperation{creerOctet}{}{\tadParams{Octet}}
    \tadOperation{ajouter}{\tadDeuxParams{Octet}{Bit}}{\tadParams{Octet}}
    \tadOperationAvecPreconditions{supprimer}{\tadDeuxParams{Octet}{Bit}}{\tadParams{Octet}}
    \tadOperation{estOctet}{\tadParams{Octet}}{\tadParams{Booleen}}
    \tadOperation{longueur}{\tadParams{Octet}}{\tadParams{\naturel}}
    \tadOperationAvecPreconditions{encoder}{\tadDeuxParams{Octet}{Code Binaire}}{\tadParams{ Chaine de caracteres}}
    \tadOperationAvecPreconditions{decoder}{\tadDeuxParams{Chaine de caracteres}{Code Binaire}}{\tadParams{Octet}}
    \tadOperationAvecPreconditions{obtenirOctet}{\tadParams{Chaine de caracteres}}{\tadParams{Octet}}
\end{tadOperations}
    
\end{tad}
    
\section{TAD Statistiques}
\begin{tad}
\tadNom{Statistiques}
\tadDependances{Caractère, Chaine de caracteres, Naturel}

\begin{tadOperations}{statistiques}

\tadOperation{statistiques}{\tadDeuxParams{{\caractere}{\naturel}}}{\tadParams{ Statistiques}}
\tadOperation{obtenirOccurences}{\tadParams{Statistiques}}{\tadParams{\naturel}}
\tadOperation{obtenirCaractere}{\tadParams{Statistiques}}{\tadParams{\caractere}}
\end{tadOperations}
\begin{tadSemantiques}{}
     statistiques : créer une statistique à partir d'un caractère et de son nombre d'occurences. \newline
     {obtenirOccurences} : Renvoie le nombre d'occurences d'un caractère et le stock dans la statistique    \newline
\end{tadSemantiques}

                          
\end{tad}



\section{TAD ArbreHuffman}

    \begin{tad}
        \tadNom{Arbre Huffman}
        
        \tadDependances{statistiques}
        
        \begin{tadOperations}{obtenirPonderation}
        
            \tadOperation{arbrehuffman}{Statistiques}{\tadParams{Arbre Huffman}}
            
            \tadOperationAvecPreconditions{fusionner}{\tadDeuxParams{Arbre Huffman}{Arbre Huffman}}{\tadParam{Arbre Huffman}}

            \tadOperation{obtenirPonderation}{\tadParams{Arbre Huffman}}{\tadParams{\naturel}}

            \tadOperationAvecPreconditions{fixerFilsDroit}{\tadDeuxParams{\paramEntree{arbred : Arbre Huffman}\paramEntreeSortie{racine : Arbre Huffman}}}

            \tadOperationAvecPreconditions{fixerFilsGauche}{\tadDeuxParams{\paramEntree{arbreg : Arbre Huffman}\paramEntreeSortie{racine : Arbre Huffman}}}

            \tadOperationAvecPreconditions{obtenirFilsGauche}{\tadParams{Arbre Huffman}}{\tadParams{Arbre Huffman}}

            \tadOperationAvecPreconditions{oHbtenirFilsDroit}{\tadParams{Arbre Huffman}}{\tadParams{Arbre Huffman}}

            \tadOperation{estFeuille}{\tadParams{Arbre Huffman}}{\tadParams{\booleen}}

            \tadOperation{estVide}{\tadParams{Arbre Huffman}}{\tadParams{\booleen}}
        \end{tadOperations}
        
        \begin{tadSemantiques}{}
        
            
            
        \end{tadSemantiques}
        
        \begin{tadPreconditions}{obtenirFilsGauche(A)}
            \tadPrecondition{fusionner(A)}{estVide(A)}
            \tadPrecondition{fixerFilsDroit(r,ad)}{non(estVide(r))}
            \tadPrecondition{fixerFilsGauche(r,ag)}{non(estVide(r))}
            \tadPrecondition{obtenirFilsGauche(A)}{non estUneFeuille(A)}
            \tadPrecondition{obtenirFilsDroit(A)}{non estUneFeuille(A)}
        \end{tadPreconditions}
           
        \end{tad}

    


\section{FileDePriorité}
\begin{tad}

\tadNom{FileDePriorité} 

\tadDependances{Booleen} 

\begin{tadOperations}{fileDePriorite}
\tadOperation{fileDePriorite}{}{\tadParams{FileDePriorite}}
\tadOperation{enfiler}{\tadDeuxParams{FileDePriorite}{Element}}{\tadParams{FileDePriorite}}
\tadOperation{defiler}{\tadParam{FileDePriorite}}{\tadParams{FileDePriorite}}
\tadOperation{longueur}{\tadParams{FiledePriorite}}{\tadParams{\naturel}}
\tadOperation{nouvelleFile}{\tadParams{FiledePriorite}}{\tadParams{FiledePriorite}}
\tadOperation{estVide}{\tadParams{FileDePriorite}}{\tadParams{\booleen}}
\tadOperation{estPresent}{\tadDeuxParams{FileDePriorite}{Element}}{\tadParams{\booleen}}
\tadOperation{obtenirElement}{\tadParams{FileDePriorite}}{\tadParams{Element}}
\end{tadOperations}

\begin{tadSemantiques}{fileDePriorite}
    \tadSemantiques{filedePriorite}{Créer une file de priotité vide}
    \tadSemantiques{enfiler}{Enfile un élément en fin de file}
    \tadSemantiques{defiler}{Défile le dernier élément de la file}
    \tadSemantiques{longueur}{Calcule la longueur de la file}
    \tadSemantiques{nouvelleFile}{Créer la nouvelle file de priorité à partir de l'actuelle}
    \tadSemantiques{estVide}{Cherche si la file de priotité est vide ou non}
    \tadSemantiques{estPresent}{Cherche si un élément donné est présent dans la file ou non}
    \tadSemantiques{obtenirElement}{Donne le dernier élément de la file}
\end{tadSemantiques}
    
\begin{tadPreconditions}
    formationNouvelArbrePondéré : non estVide(defiler(defiler(fdp)))\newline
\end{tadPreconditions}

\end{tad}
\section{TAD TableDeCodage}

\begin{tad}
\tadNom{TableDeCodage}
\tadDependances{ArbreHuffmann}

\begin{tadOperations}{obtenirTableDeCodage}
    \tadOperation{obtenirTableDeCodage}{\tadParams{ArbreHuffmann}}{\tadParams{TableDeCodage}}
\tadOperation{}{}{}
\end{tadOperations}

\begin{tadSemantiques}{}
    obtenirTableDeCodage : Renvoie la table de codage grâce à un arbre de Huffmann \newline
\end{tadSemantiques}
    
\begin{tadPreconditions}{obtenirTableDeCodage}
    obtenirTableDeCodage : non estVide(arbrebinaire)
\end{tadPreconditions}
\end{tad}
\end{document}

